\documentclass[11pt]{article}

\usepackage[utf8]{inputenc}
\usepackage[T1]{fontenc}
\usepackage{lmodern}
\usepackage{geometry}
\usepackage{amsmath,amssymb,amsfonts,amsthm,mathtools}
\usepackage{bm}
\usepackage{enumitem}
\usepackage{microtype}
\usepackage{hyperref}
\usepackage{titlesec}
\usepackage{booktabs}
\usepackage{array}
\usepackage{setspace}
\usepackage{mathrsfs}
\usepackage{physics}

\geometry{margin=1in}
\hypersetup{colorlinks=true, linkcolor=black, urlcolor=blue, citecolor=black}
\setlist[enumerate]{topsep=0.4em,itemsep=0.35em}
\setlist[itemize]{topsep=0.3em,itemsep=0.25em}
\titleformat{\section}{\large\bfseries}{\thesection.}{0.5em}{}
\titleformat{\subsection}{\normalsize\bfseries}{\thesubsection.}{0.5em}{}
\setstretch{1.06}

\newcommand{\Aether}{\AE{}ther}
\newcommand{\Aflow}{\AE{}ther-flow}
\newcommand{\TheoryName}{\Aflow{} Interpretation of Relativity}
\newcommand{\TheoryProgram}{\Aether{} / \Aflow{} framework}
\newcommand{\Sub}{\mathcal A}
\newcommand{\ObserverMap}{\Xi}
\newcommand{\BridgeMetric}{\mathfrak g}
\newcommand{\ObserverMetric}{g^{\mathrm{br}}}
\newcommand{\CandidateMetric}{g^{\mathrm{cand}}}
\newcommand{\CompletedMetric}{g^{\sharp}}
\newcommand{\BaseShift}{\Sigma^{\mathrm{IER}}}
\newcommand{\ResponseShift}{\Sigma^{\mathrm{resp}}}
\newcommand{\CompShift}{\Sigma^{\mathrm{cmp}}}
\newcommand{\BaseLift}{\widehat\Sigma^{\mathrm{IER}}}
\newcommand{\ResponseLift}{\widehat\Sigma^{\mathrm{resp}}}
\newcommand{\CompLift}{\widehat\Sigma^{\mathrm{cmp}}}
\newcommand{\ReLongFour}{\widetilde{\mathcal R}_{\mu\nu}^{(\chi,\ge 4)}}
\newcommand{\ResidualCore}{\mathcal V_{\mu\nu}^{\mathrm{res}}}
\newcommand{\ResidualLift}{\widehat{\mathcal V}^{\mathrm{res}}}
\newcommand{\SubEinLin}{\widehat{\mathcal E}}
\newcommand{\SubGaugeVec}{\widehat{\mathcal B}}
\newcommand{\SubGaugeBook}{\widehat{\mathcal G}}
\newcommand{\SubReducedEin}{\widehat{\mathcal L}}
\newcommand{\FreeProj}{\Pi_{\mathrm{free}}}
\newcommand{\GaugeAux}{Y}
\newcommand{\ReducedEin}{\mathcal L_{\ObserverMetric}}
\newcommand{\GaugeBook}{\mathcal G_{\ObserverMetric}}
\newcommand{\EinLin}{\mathcal E_{\ObserverMetric}}
\newcommand{\EinQuad}{\mathcal Q_{\ge 2}^{\mathrm{Ein}}}
\newcommand{\CompClass}{\mathscr C_{\mathrm{cmp}}}
\newcommand{\CompFree}{\mathscr C_{\mathrm{cmp}}^{\mathrm{free}}}
\newcommand{\CompForced}{\mathscr C_{\mathrm{cmp}}^{\mathrm{for}}}
\newcommand{\CompKernel}{\mathscr K_{\mathrm{cmp}}}
\newcommand{\MicFunctional}{\mathscr F_{\mathrm{mic}}}
\newcommand{\PrimFunctional}{\mathscr F_{\mathrm{prim}}}
\newcommand{\CarrierFunctional}{\mathscr F_{\mathrm{car}}}
\newcommand{\CarrierToPrim}{\mathscr F_{\mathrm{car}\to\mathrm{prim}}}
\newcommand{\DeepFunctional}{\mathscr F_{\mathrm{deep}}}
\newcommand{\DeepToCar}{\mathscr F_{\mathrm{deep}\to\mathrm{car}}}
\newcommand{\ProtoFunctional}{\mathscr F_{\mathrm{proto}}}
\newcommand{\ProtoToDeep}{\mathscr F_{\mathrm{proto}\to\mathrm{deep}}}
\newcommand{\FreeStiff}{\kappa_{\mathrm{free}}}
\newcommand{\PrimStrain}{K}
\newcommand{\PrimNeutral}{J}
\newcommand{\PrimFree}{F}
\newcommand{\CarrierClass}{\mathscr H_{\mathrm{car}}}
\newcommand{\CarrierProj}{\Pi_{\mathrm{str}}}
\newcommand{\CarrierPerp}{\Pi_{\mathrm{str}}^{\perp}}
\newcommand{\CarrierGap}{\kappa_{\mathrm{str}}}
\newcommand{\CarrierSeed}{\mathfrak S}
\newcommand{\RelabelSeed}{\mathfrak R}
\newcommand{\FreeSeed}{\mathfrak F}
\newcommand{\PreCarGap}{\kappa_{\mathrm{car},0}}
\newcommand{\CurrentGap}{\kappa_{\mathrm{cur}}}
\newcommand{\HomGap}{\kappa_{\mathrm{hom}}}
\newcommand{\StrainMult}{\Lambda}
\newcommand{\NeutralMult}{\Omega}
\newcommand{\FreeMult}{\Theta}
\newcommand{\epsIR}{\varepsilon}
\newcommand{\DeepTensorClass}{\mathscr X_{\mathrm{car}}}
\newcommand{\DeepRelabelClass}{\mathscr X_{\mathrm{cur}}}
\newcommand{\DeepFreeClass}{\mathscr X_{\mathrm{hom}}}
\newcommand{\DeepTensor}{\mathfrak X}
\newcommand{\DeepRelabel}{\mathfrak J}
\newcommand{\DeepFree}{\mathfrak W}
\newcommand{\CarRead}{\mathcal R_{\mathrm{car}}}
\newcommand{\CurRead}{\mathcal R_{\mathrm{cur}}}
\newcommand{\HomRead}{\mathcal R_{\mathrm{hom}}}
\newcommand{\ProtoTensorClass}{\mathscr Y_{\mathrm{car}}}
\newcommand{\ProtoRelabelClass}{\mathscr Y_{\mathrm{cur}}}
\newcommand{\ProtoFreeClass}{\mathscr Y_{\mathrm{hom}}}
\newcommand{\ProtoTensor}{\mathfrak Y}
\newcommand{\ProtoRelabel}{\mathfrak K}
\newcommand{\ProtoFree}{\mathfrak Z}
\newcommand{\ProtoCarProj}{\Pi_{\mathrm{deep,car}}}
\newcommand{\ProtoCurProj}{\Pi_{\mathrm{deep,cur}}}
\newcommand{\ProtoHomProj}{\Pi_{\mathrm{deep,hom}}}
\newcommand{\ProtoCarPerp}{\Pi_{\mathrm{deep,car}}^{\perp}}
\newcommand{\ProtoCurPerp}{\Pi_{\mathrm{deep,cur}}^{\perp}}
\newcommand{\ProtoHomPerp}{\Pi_{\mathrm{deep,hom}}^{\perp}}
\newcommand{\ProtoCarRead}{\widetilde{\mathcal R}_{\mathrm{car}}}
\newcommand{\ProtoCurRead}{\widetilde{\mathcal R}_{\mathrm{cur}}}
\newcommand{\ProtoHomRead}{\widetilde{\mathcal R}_{\mathrm{hom}}}
\newcommand{\ProtoCarGap}{\gamma_{\mathrm{deep,car}}}
\newcommand{\ProtoCurGap}{\gamma_{\mathrm{deep,cur}}}
\newcommand{\ProtoHomGap}{\gamma_{\mathrm{deep,hom}}}
\newcommand{\UltraFunctional}{\mathscr F_{\mathrm{micro}}}
\newcommand{\UltraToProto}{\mathscr F_{\mathrm{micro}\to\mathrm{proto}}}
\newcommand{\UltraTensorClass}{\mathscr Z_{\mathrm{car}}}
\newcommand{\UltraRelabelClass}{\mathscr Z_{\mathrm{cur}}}
\newcommand{\UltraFreeClass}{\mathscr Z_{\mathrm{hom}}}
\newcommand{\UltraTensor}{\mathfrak M}
\newcommand{\UltraRelabel}{\mathfrak N}
\newcommand{\UltraFree}{\mathfrak P}
\newcommand{\UltraCarProj}{\Pi_{\mathrm{mic,car}}}
\newcommand{\UltraCurProj}{\Pi_{\mathrm{mic,cur}}}
\newcommand{\UltraHomProj}{\Pi_{\mathrm{mic,hom}}}
\newcommand{\UltraCarPerp}{\Pi_{\mathrm{mic,car}}^{\perp}}
\newcommand{\UltraCurPerp}{\Pi_{\mathrm{mic,cur}}^{\perp}}
\newcommand{\UltraHomPerp}{\Pi_{\mathrm{mic,hom}}^{\perp}}
\newcommand{\UltraCarGap}{\delta_{\mathrm{mic,car}}}
\newcommand{\UltraCurGap}{\delta_{\mathrm{mic,cur}}}
\newcommand{\UltraHomGap}{\delta_{\mathrm{mic,hom}}}
\newcommand{\TensorDefect}{\mathcal D_{\mathrm{car}}}
\newcommand{\CurDefect}{\mathcal D_{\mathrm{cur}}}
\newcommand{\HomDefect}{\mathcal D_{\mathrm{hom}}}
\newcommand{\UltraCarRead}{\mathcal S_{\mathrm{car}}}
\newcommand{\UltraCurRead}{\mathcal S_{\mathrm{cur}}}
\newcommand{\UltraHomRead}{\mathcal S_{\mathrm{hom}}}

\newtheorem{definition}{Definition}
\newtheorem{proposition}{Proposition}
\newtheorem{remark}{Remark}
\newtheorem{corollary}{Corollary}
\newtheorem{theorem}{Theorem}

\title{\textbf{The \TheoryName{}}\\[0.4em]
\large Substrate-Response / Self-Response Charge-Polarization Candidate\\
\large Quartic-Completion Selection-Principle Primitive-Reservoir Carrier-Origin\\
\large Precursor-Selection Exact-Preservation Normal-Form Microphysical-Origin Note}
\author{}
\date{}

\begin{document}

\maketitle

\begin{abstract}
The exact-preservation normal-form origin note already shows that the raw deeper readout spaces are not primitive at present scope. They are the projected unsuppressed sectors of ambient substrate spaces
\[
\ProtoTensorClass,\qquad
\ProtoRelabelClass,\qquad
\ProtoFreeClass
\]
with orthogonal projectors
\[
\ProtoCarProj:\ProtoTensorClass\to\DeepTensorClass,
\qquad
\ProtoCurProj:\ProtoRelabelClass\to\DeepRelabelClass,
\qquad
\ProtoHomProj:\ProtoFreeClass\to\DeepFreeClass,
\]
factorized visible maps into
\[
\CarrierClass,\qquad
T^*\Sub,\qquad
\CompFree,
\]
and positively gapped ambient spectators. What remained open is one layer deeper again: why those ambient substrate spaces, those ambient projectors, those factorized visible maps, and those positive spectator gaps should exist at all rather than being the present deepest disciplined postulate.

The present note gives that next origin step at disciplined microphysical exact-preservation scope. It introduces still more primitive microphysical spaces
\[
\UltraTensorClass,\qquad
\UltraRelabelClass,\qquad
\UltraFreeClass
\]
equipped with closed defect operators
\[
\TensorDefect:\UltraTensorClass\to \mathscr D_{\mathrm{car}},
\qquad
\CurDefect:\UltraRelabelClass\to \mathscr D_{\mathrm{cur}},
\qquad
\HomDefect:\UltraFreeClass\to \mathscr D_{\mathrm{hom}},
\]
whose exact-preserving kernel sectors are precisely
\[
\ProtoTensorClass=\ker\TensorDefect,\qquad
\ProtoRelabelClass=\ker\CurDefect,\qquad
\ProtoFreeClass=\ker\HomDefect.
\]
The associated orthogonal kernel projectors
\[
\UltraCarProj:\UltraTensorClass\to\ProtoTensorClass,
\qquad
\UltraCurProj:\UltraRelabelClass\to\ProtoRelabelClass,
\qquad
\UltraHomProj:\UltraFreeClass\to\ProtoFreeClass,
\]
are forced by Hilbert orthogonality once the exact-preserving sectors are identified as the defect-null directions. Coercive defect estimates
\[
\norm{\TensorDefect \UltraTensor}_{\mathrm{def,car}}^2
\ge \UltraCarGap\norm{\UltraCarPerp\UltraTensor}_{\mathrm{micro,car}}^2,
\qquad
\norm{\CurDefect \UltraRelabel}_{\mathrm{def,cur}}^2
\ge \UltraCurGap\norm{\UltraCurPerp\UltraRelabel}_{\mathrm{micro,cur}}^2,
\]
\[
\norm{\HomDefect \UltraFree}_{\mathrm{def,hom}}^2
\ge \UltraHomGap\norm{\UltraHomPerp\UltraFree}_{\mathrm{micro,hom}}^2,
\qquad
\UltraCarGap,\UltraCurGap,\UltraHomGap>0,
\]
make the defect-bearing complements strictly positive-gap spectators. The visible microphysical maps are then forced to factor through the exact-preserving kernel sectors and through the ambient raw-sector projectors as
\[
\UltraCarRead=\CarRead\circ\ProtoCarProj\circ\UltraCarProj,
\qquad
\UltraCurRead=\CurRead\circ\ProtoCurProj\circ\UltraCurProj,
\qquad
\UltraHomRead=\HomRead\circ\ProtoHomProj\circ\UltraHomProj,
\]
because the fixed carrier-origin functional still has exactly three visible slots and no others:
\[
\CarrierClass,\qquad
T^*\Sub,\qquad
\CompFree.
\]
The defect-bearing microphysical complements therefore become the still deeper exact-preserving spectators. Inside the recovered ambient exact-preserving sectors, the ambient complements remain the same positively gapped spectators of the previous note, and inside the projected raw spaces the hidden readout-null directions remain the same positively gapped kernel sectors recorded there.

Eliminating the microphysical defect layer reproduces exactly the same ambient exact-preservation functional
\[
\ProtoFunctional[\ProtoTensor,\ProtoRelabel,\ProtoFree]
\]
already recorded one layer above. The previously recorded ambient elimination then reproduces exactly the same exact-preservation normal form
\[
\DeepFunctional[\DeepTensor,\DeepRelabel,\DeepFree],
\]
and the already recorded downstream eliminations again reproduce exactly the same carrier-origin functional, the same primitive reservoir, the same microscopic-equilibrium reservoir, the same \(\Pi_{\mathrm{free}}H=0\) outcome, the same substrate and observer completion solves, and the same one operative metric. In particular, the same carrier-origin functional yields the same primitive reservoir, the same microscopic-equilibrium completion reservoir, the same completion solve \(\Sigma^{\mathrm{cmp}}\), and the same observer-equation closure.

The scope remains disciplined. This note does not claim a final ultraviolet completion or a uniqueness theorem for every conceivable more primitive substrate sector. Its positive claim is narrower and exact: at the current quadratic exact-preservation scope, the ambient projected-sector construction itself is recovered as the defect-null exact-preserving sector of a still more primitive microphysical layer with coercively gapped defect directions, while preserving the same downstream one-metric completion package.
\end{abstract}

\section{Introduction}

The active charge-polarization line has now crossed the candidate, gate, controlled-limit, residual-sector, redesign, substrate-anchoring, substrate-action, kernel-origin, microscopic-equilibrium deeper-origin, equilibrium-reservoir primitive-origin, primitive-reservoir carrier-origin, carrier-origin precursor-selection, carrier-origin precursor-selection deeper-origin, and exact-preservation normal-form origin steps on the same one-metric GR route. The immediate burden therefore narrows again. It is not another packaging pass. It is not another same-layer rewrite of the preceding note. It is not stronger promotion language. It is not, by default, another anchored positive-pair continuation. It is to go one layer below the ambient projected-sector construction now on record and explain why that construction itself arises from still more primitive substrate structure.

That burden has five tightly linked parts. First, one must explain why the ambient substrate spaces
\[
\ProtoTensorClass,\qquad
\ProtoRelabelClass,\qquad
\ProtoFreeClass
\]
are forced at all. Second, one must explain why the ambient projectors
\[
\ProtoCarProj,\qquad
\ProtoCurProj,\qquad
\ProtoHomProj
\]
onto the raw deeper sectors are themselves honest induced structure rather than another same-layer input. Third, one must explain why the only visible maps from the still deeper layer are the codomain-preserving maps into
\[
\CarrierClass,\qquad
T^*\Sub,\qquad
\CompFree.
\]
Fourth, one must explain why the hidden exact-preserving and exact-breaking directions carry coercive quadratic sectors and why those hidden directions are positively gapped. Fifth, one must re-show that eliminating this still more primitive layer returns exactly the same downstream package already recorded above it: the same carrier-origin functional, the same primitive reservoir, the same microscopic-equilibrium reservoir, the same \(\Sigma^{\mathrm{cmp}}\) solve, the same \(\Pi_{\mathrm{free}}H=0\) verdict, the same observer-equation closure, and the same one operative metric.

The present note addresses exactly those tasks. Its key move is to step one layer below the ambient projected-sector construction itself and to introduce a still more primitive microphysical exact-preservation layer. In that microphysical layer, the admissible exact-preserving ambient sector is defined as the common defect-null sector of three closed defect operators. Their orthogonal kernel projectors then recover the ambient spaces as honest induced subspaces rather than unexplained primitives. Coercive defect estimates make the exact-breaking complements strictly positive-gap spectators. Once those defect-bearing complements are eliminated, the same ambient projected-sector construction of the previous note is recovered exactly. The ambient elimination then again recovers the same raw deeper readout sector, and from there the same downstream carrier-origin, primitive-reservoir, microscopic-equilibrium, completion, and observer-equation package follows without modification.

\paragraph{Framework and claim boundary.}
\TheoryProgram{} denotes the broader \Aether{} / \Aflow{} framework built on the claims that the \Aether{} is the underlying four-dimensional substrate of reality, the \Aflow{} is its intrinsic ordered motion, observed three-dimensional space is the local experiential slice of that deeper substrate, S-time is the experienced order of change arising from matter, light, and the \Aflow{}, and observed expansion is the three-dimensional appearance of deeper four-dimensional ordered motion. Within that framework, \TheoryName{} denotes the exact-closure theory in which the effective gravitational dynamics are adopted to be exactly Einsteinian with universal matter coupling. In the active sequence, \emph{adoption} means use of that established relativistic dynamics without claiming substrate derivation, while \emph{derivation} is reserved for a first-principles recovery from explicit substrate variables.

The present note stays strictly inside that boundary. It does not claim that every possible still more primitive substrate sector has now been classified. It does not widen the current theorem boundary. It does make one narrower positive move: it shows that the ambient projected-sector construction of the previous note is itself recovered as the defect-null exact-preserving sector of a still more primitive microphysical layer with coercively gapped defect directions, while preserving the same carrier-origin functional, the same primitive reservoir, the same microscopic-equilibrium reservoir, the same \(\Sigma^{\mathrm{cmp}}\) solve, and the same one operative metric.

\section{Fixed Inputs from the Recorded Completion Line}

\begin{definition}[Fixed charge-polarization branch and source]
\label{def:fixed}
The present note takes as fixed the same charge-polarization branch
\begin{equation}
\CandidateMetric
=
\ObserverMetric+\BaseShift+\ResponseShift,
\qquad
\CompletedMetric
=
\ObserverMetric+\BaseShift+\ResponseShift+\CompShift,
\label{eq:metrics}
\end{equation}
with lifted completion variables \(\BaseLift,\ResponseLift,\CompLift\) satisfying
\begin{equation}
\ObserverMap^*\BaseLift=\BaseShift,
\qquad
\ObserverMap^*\ResponseLift=\ResponseShift,
\qquad
\ObserverMap^*\CompLift=\CompShift.
\label{eq:lifts}
\end{equation}
The unresolved observer-side quartic tensor and its fixed substrate lift are
\begin{equation}
\ResidualCore
\equiv
\ReLongFour
+\EinQuad[\BaseShift]
+\EinLin(\ResponseShift),
\label{eq:vres}
\end{equation}
\begin{equation}
\ResidualLift_{AB}
\equiv
\widehat{\mathcal R}_{AB}^{(\chi,\ge 4)}
+\widehat{\mathcal Q}_{AB}^{\mathrm{Ein}}[\BaseLift]
+\widehat{\mathcal E}_{AB}(\ResponseLift),
\qquad
\ObserverMap^*\ResidualLift=\ResidualCore.
\label{eq:liftedsource}
\end{equation}
The fixed orders are
\begin{equation}
\BaseShift=\mathcal O_{\ge 2},
\qquad
\ResponseShift=\mathcal O_{\ge 4},
\qquad
\CompShift=\mathcal O_{\ge 4},
\qquad
\ResidualLift=\mathcal O_{\ge 4},
\label{eq:orders}
\end{equation}
and on every controlled infrared family,
\begin{equation}
\BaseShift=\mathcal O(\epsIR^2),
\qquad
\ResponseShift=\mathcal O(\epsIR^4),
\qquad
\CompShift=\mathcal O(\epsIR^4),
\qquad
\ResidualLift=\mathcal O(\epsIR^4).
\label{eq:irorders}
\end{equation}
\end{definition}

\begin{definition}[Fixed bridge-response operators]
\label{def:operators}
Let \(\widehat\nabla\) denote the Levi-Civita connection of the unshifted bridge metric \(\BridgeMetric\). For any observer-compatible symmetric substrate tensor \(H_{AB}\), define
\begin{equation}
\SubGaugeVec(H)_A
\equiv
\widehat\nabla^B\!\left(
H_{AB}
-\frac12 \BridgeMetric_{AB}\BridgeMetric^{CD}H_{CD}
\right),
\label{eq:subB}
\end{equation}
\begin{equation}
\SubGaugeBook(H)_{AB}
\equiv
\widehat\nabla_{(A}\SubGaugeVec(H)_{B)}
-\frac12 \BridgeMetric_{AB}\widehat\nabla^C\SubGaugeVec(H)_C,
\label{eq:subG}
\end{equation}
\begin{equation}
\SubReducedEin(H)_{AB}
\equiv
\SubEinLin(H)_{AB}-\SubGaugeBook(H)_{AB},
\label{eq:subL}
\end{equation}
chosen so that
\begin{equation}
\ObserverMap^*\bigl(\SubEinLin(H)\bigr)=\EinLin(\ObserverMap^*H),
\qquad
\ObserverMap^*\bigl(\SubGaugeBook(H)\bigr)=\GaugeBook(\ObserverMap^*H),
\label{eq:intertwine1}
\end{equation}
and hence
\begin{equation}
\ObserverMap^*\bigl(\SubReducedEin(H)\bigr)
=
\ReducedEin(\ObserverMap^*H).
\label{eq:intertwine2}
\end{equation}
\end{definition}

\begin{definition}[Admissible completion class and free sector]
\label{def:class}
Let \(\CompClass\) denote the same observer-compatible reduced-harmonic Coulomb-plus-regular completion class used in the redesign, anchoring, action, kernel-origin, microscopic-equilibrium deeper-origin, equilibrium-reservoir primitive-origin, primitive-reservoir carrier-origin, carrier-origin precursor-selection, and precursor-selection deeper-origin notes. The free homogeneous sector is
\begin{equation}
\CompFree
\equiv
\left\{
H\in\CompClass:\SubReducedEin(H)=0
\right\},
\label{eq:freeclass}
\end{equation}
and \(\CompForced\) denotes its orthogonal complement inside \(\CompClass\) with respect to the bridge-metric \(L^2\) pairing
\begin{equation}
\langle H_1,H_2\rangle_{\mathrm{cmp}}
\equiv
\int_{\Sub} d^4X\,\sqrt{G}\,
\BridgeMetric^{AC}\BridgeMetric^{BD}(H_1)_{AB}(H_2)_{CD}.
\label{eq:inner}
\end{equation}
Let \(\FreeProj:\CompClass\to\CompFree\) be the orthogonal projector. The fixed source satisfies
\begin{equation}
\ResidualLift\in\CompForced,
\qquad
\FreeProj\ResidualLift=0.
\label{eq:sourcefree}
\end{equation}
\end{definition}

\begin{remark}
The present note does not alter the source \(\ResidualLift\), the reduced operator \(\SubReducedEin\), the one-metric matter route, or the admissible completion class. The new burden is one layer deeper than the previous note: why the ambient projected-sector construction itself should arise from more primitive substrate structure.
\end{remark}

\section{Fixed Ambient Exact-Preservation Target}

\begin{definition}[Carrier-origin block fixed by the previous notes]
\label{def:carrier}
The primitive-reservoir carrier-origin note fixed a still-deeper carrier block with tensor seed
\[
\CarrierSeed\in\CarrierClass,
\qquad
\CarrierClass=\CompClass\oplus \CompClass^\perp,
\]
orthogonal projector
\[
\CarrierProj:\CarrierClass\to\CompClass,
\qquad
\CarrierPerp\equiv I-\CarrierProj,
\]
relabeling-current seed \(\RelabelSeed_A\in T^*\Sub\), and source-free homogeneous seed \(\FreeSeed_{AB}\in\CompFree\). Its functional is
\begin{equation}
\begin{aligned}
\CarrierFunctional[\CarrierSeed,\RelabelSeed,\FreeSeed]
\equiv
\int_{\Sub} d^4X\,\sqrt{G}\,
\Bigl[
&\frac12 (\CarrierProj\CarrierSeed)^{AB}
\SubEinLin(\CarrierProj\CarrierSeed)_{AB}
+\frac{\CarrierGap}{2}(\CarrierPerp\CarrierSeed)^{AB}(\CarrierPerp\CarrierSeed)_{AB}
\\
&+\frac12\bigl(\RelabelSeed_A-\SubGaugeVec(\CarrierProj\CarrierSeed)_A\bigr)
\bigl(\RelabelSeed^A-\SubGaugeVec(\CarrierProj\CarrierSeed)^A\bigr)
\\
&+\frac{\FreeStiff}{2}\FreeSeed^{AB}\FreeSeed_{AB}
+(\ResidualLift)^{AB}(\CarrierProj\CarrierSeed)_{AB}
\Bigr],
\\
&\qquad\qquad \CarrierGap>0,\qquad \FreeStiff>0.
\end{aligned}
\label{eq:car}
\end{equation}
The primitive readout from that carrier block is
\begin{equation}
\PrimStrain\equiv \CarrierProj\CarrierSeed,
\qquad
\PrimNeutral\equiv \RelabelSeed,
\qquad
\PrimFree\equiv \FreeSeed.
\label{eq:readout}
\end{equation}
\end{definition}

\begin{definition}[Exact-preservation normal-form target]
\label{def:deep}
The precursor-selection deeper-origin note fixed raw deeper readout spaces
\[
\DeepTensorClass,\qquad
\DeepRelabelClass,\qquad
\DeepFreeClass
\]
and bounded surjective linear maps
\begin{equation}
\CarRead:\DeepTensorClass\to\CarrierClass,
\qquad
\CurRead:\DeepRelabelClass\to T^*\Sub,
\qquad
\HomRead:\DeepFreeClass\to\CompFree.
\label{eq:maps}
\end{equation}
For every
\[
\DeepTensor\in\DeepTensorClass,\qquad
\DeepRelabel\in\DeepRelabelClass,\qquad
\DeepFree\in\DeepFreeClass,
\]
write the orthogonal decompositions
\begin{equation}
\DeepTensor=\DeepTensor^{\parallel}+\DeepTensor^{\perp},
\qquad
\DeepTensor^{\parallel}\in(\ker\CarRead)^\perp,
\qquad
\DeepTensor^{\perp}\in\ker\CarRead,
\label{eq:Xsplit}
\end{equation}
\begin{equation}
\DeepRelabel=\DeepRelabel^{\parallel}+\DeepRelabel^{\perp},
\qquad
\DeepRelabel^{\parallel}\in(\ker\CurRead)^\perp,
\qquad
\DeepRelabel^{\perp}\in\ker\CurRead,
\label{eq:Jsplit}
\end{equation}
\begin{equation}
\DeepFree=\DeepFree^{\parallel}+\DeepFree^{\perp},
\qquad
\DeepFree^{\parallel}\in(\ker\HomRead)^\perp,
\qquad
\DeepFree^{\perp}\in\ker\HomRead.
\label{eq:Wsplit}
\end{equation}
Its exact-preservation normal-form functional is
\begin{equation}
\begin{aligned}
\DeepFunctional[\DeepTensor,\DeepRelabel,\DeepFree]
\equiv\;
&\CarrierFunctional[\CarRead\DeepTensor,\CurRead\DeepRelabel,\HomRead\DeepFree]
\\
&+\frac{\PreCarGap}{2}\norm{\DeepTensor^{\perp}}_{\mathrm{deep,car}}^2
+\frac{\CurrentGap}{2}\norm{\DeepRelabel^{\perp}}_{\mathrm{deep,cur}}^2
+\frac{\HomGap}{2}\norm{\DeepFree^{\perp}}_{\mathrm{deep,hom}}^2,
\\
&\qquad\qquad
\PreCarGap>0,\qquad
\CurrentGap>0,\qquad
\HomGap>0.
\end{aligned}
\label{eq:Fdeep}
\end{equation}
\end{definition}

\begin{definition}[Recorded ambient projected-sector construction]
\label{def:proto}
The exact-preservation normal-form origin note fixed ambient substrate spaces
\[
\ProtoTensorClass,\qquad
\ProtoRelabelClass,\qquad
\ProtoFreeClass
\]
containing closed projected sectors identified respectively with
\[
\DeepTensorClass,\qquad
\DeepRelabelClass,\qquad
\DeepFreeClass.
\]
Let
\begin{equation}
\ProtoCarProj:\ProtoTensorClass\to\DeepTensorClass,
\qquad
\ProtoCurProj:\ProtoRelabelClass\to\DeepRelabelClass,
\qquad
\ProtoHomProj:\ProtoFreeClass\to\DeepFreeClass
\label{eq:proto_proj}
\end{equation}
denote the orthogonal projectors onto those closed subspaces, and set
\begin{equation}
\ProtoCarPerp\equiv I-\ProtoCarProj,
\qquad
\ProtoCurPerp\equiv I-\ProtoCurProj,
\qquad
\ProtoHomPerp\equiv I-\ProtoHomProj.
\label{eq:proto_perp}
\end{equation}
Define the ambient visible readout maps by
\begin{equation}
\ProtoCarRead\equiv \CarRead\circ\ProtoCarProj,
\qquad
\ProtoCurRead\equiv \CurRead\circ\ProtoCurProj,
\qquad
\ProtoHomRead\equiv \HomRead\circ\ProtoHomProj.
\label{eq:proto_maps}
\end{equation}
Then
\begin{equation}
\ProtoCarRead\,\ProtoCarPerp=0,
\qquad
\ProtoCurRead\,\ProtoCurPerp=0,
\qquad
\ProtoHomRead\,\ProtoHomPerp=0,
\label{eq:proto_kill}
\end{equation}
and the ambient visible maps have the same codomains
\[
\CarrierClass,\qquad
T^*\Sub,\qquad
\CompFree.
\]
\end{definition}

\begin{definition}[Recorded ambient exact-preservation functional]
\label{def:Fproto}
The previously recorded ambient exact-preservation functional is
\begin{equation}
\begin{aligned}
\ProtoFunctional[\ProtoTensor,\ProtoRelabel,\ProtoFree]
\equiv\;
&\DeepFunctional[
\ProtoCarProj\ProtoTensor,\ProtoCurProj\ProtoRelabel,
\\
&\hspace{1.7em}\ProtoHomProj\ProtoFree]
\\
&+\frac{\ProtoCarGap}{2}\norm{\ProtoCarPerp\ProtoTensor}_{\mathrm{proto,car}}^2
+\frac{\ProtoCurGap}{2}\norm{\ProtoCurPerp\ProtoRelabel}_{\mathrm{proto,cur}}^2
+\frac{\ProtoHomGap}{2}\norm{\ProtoHomPerp\ProtoFree}_{\mathrm{proto,hom}}^2,
\\
&\qquad\qquad
\ProtoCarGap>0,\qquad
\ProtoCurGap>0,\qquad
\ProtoHomGap>0.
\end{aligned}
\label{eq:Fproto}
\end{equation}
\end{definition}

\begin{definition}[Recorded downstream reservoirs]
\label{def:downstream}
The same previous notes already fixed the primitive reservoir
\begin{equation}
\begin{aligned}
\PrimFunctional[\PrimStrain,\PrimNeutral,\PrimFree]
\equiv
\int_{\Sub} d^4X\,\sqrt{G}\,
\Bigl[
&\frac12 \PrimStrain^{AB}\SubEinLin(\PrimStrain)_{AB}
+\frac12\bigl(\PrimNeutral_A-\SubGaugeVec(\PrimStrain)_A\bigr)
\bigl(\PrimNeutral^A-\SubGaugeVec(\PrimStrain)^A\bigr)
\\
&+\frac{\FreeStiff}{2}\PrimFree^{AB}\PrimFree_{AB}
+(\ResidualLift)^{AB}\PrimStrain_{AB}
\Bigr],
\end{aligned}
\label{eq:primitive}
\end{equation}
and the microscopic-equilibrium completion reservoir
\begin{equation}
\begin{aligned}
\MicFunctional[H,\GaugeAux]
=
\int_{\Sub} d^4X\,\sqrt{G}\,
\Bigl[
&\frac12 H^{AB}\SubEinLin(H)_{AB}
+\frac12\bigl(\GaugeAux_A-\SubGaugeVec(H)_A\bigr)
\bigl(\GaugeAux^A-\SubGaugeVec(H)^A\bigr)
\\
&+\frac{\FreeStiff}{2}(\FreeProj H)^{AB}(\FreeProj H)_{AB}
+(\ResidualLift)^{AB}H_{AB}
\Bigr].
\end{aligned}
\label{eq:mic}
\end{equation}
\end{definition}

\begin{remark}
The present manuscript takes the ambient projected-sector functional \eqref{eq:Fproto} as the immediate target that must now be derived or justified from a still more primitive microphysical layer. Once \eqref{eq:Fproto} is recovered, the previously recorded ambient elimination recovers the same exact-preservation normal form \eqref{eq:Fdeep}; once \eqref{eq:Fdeep} is recovered, the previously recorded elimination chain to \eqref{eq:car}, \eqref{eq:primitive}, \eqref{eq:mic}, \(\Sigma^{\mathrm{cmp}}\), \(\Pi_{\mathrm{free}}H=0\), and the one operative metric becomes available again.
\end{remark}

\section{Still-Deeper Microphysical Exact-Preservation Sector}

\begin{definition}[Microphysical defect operators and exact-preserving ambient kernels]
\label{def:ultra}
Let
\[
\UltraTensorClass,\qquad
\UltraRelabelClass,\qquad
\UltraFreeClass
\]
be Hilbert spaces equipped with closed densely defined defect operators
\begin{equation}
\TensorDefect:\UltraTensorClass\to\mathscr D_{\mathrm{car}},
\qquad
\CurDefect:\UltraRelabelClass\to\mathscr D_{\mathrm{cur}},
\qquad
\HomDefect:\UltraFreeClass\to\mathscr D_{\mathrm{hom}},
\label{eq:defects}
\end{equation}
whose kernels are closed. Define the exact-preserving ambient sectors by
\begin{equation}
\ProtoTensorClass\equiv \ker\TensorDefect,
\qquad
\ProtoRelabelClass\equiv \ker\CurDefect,
\qquad
\ProtoFreeClass\equiv \ker\HomDefect,
\label{eq:proto_kernels}
\end{equation}
and let
\begin{equation}
\UltraCarProj:\UltraTensorClass\to\ProtoTensorClass,
\qquad
\UltraCurProj:\UltraRelabelClass\to\ProtoRelabelClass,
\qquad
\UltraHomProj:\UltraFreeClass\to\ProtoFreeClass
\label{eq:ultra_proj}
\end{equation}
be the orthogonal projectors onto those kernels. Set
\begin{equation}
\UltraCarPerp\equiv I-\UltraCarProj,
\qquad
\UltraCurPerp\equiv I-\UltraCurProj,
\qquad
\UltraHomPerp\equiv I-\UltraHomProj.
\label{eq:ultra_perp}
\end{equation}
Assume the coercive defect estimates
\begin{equation}
\norm{\TensorDefect\UltraTensor}_{\mathrm{def,car}}^2
\ge
\UltraCarGap\,\norm{\UltraCarPerp\UltraTensor}_{\mathrm{micro,car}}^2,
\qquad
\UltraCarGap>0,
\label{eq:defect_car}
\end{equation}
\begin{equation}
\norm{\CurDefect\UltraRelabel}_{\mathrm{def,cur}}^2
\ge
\UltraCurGap\,\norm{\UltraCurPerp\UltraRelabel}_{\mathrm{micro,cur}}^2,
\qquad
\UltraCurGap>0,
\label{eq:defect_cur}
\end{equation}
\begin{equation}
\norm{\HomDefect\UltraFree}_{\mathrm{def,hom}}^2
\ge
\UltraHomGap\,\norm{\UltraHomPerp\UltraFree}_{\mathrm{micro,hom}}^2,
\qquad
\UltraHomGap>0.
\label{eq:defect_hom}
\end{equation}
\end{definition}

\begin{definition}[Microphysical visible maps and exact-preservation functional]
\label{def:Fultra}
Define the visible microphysical maps by
\begin{equation}
\UltraCarRead\equiv \ProtoCarRead\circ\UltraCarProj
=
\CarRead\circ\ProtoCarProj\circ\UltraCarProj,
\label{eq:ultra_car_map}
\end{equation}
\begin{equation}
\UltraCurRead\equiv \ProtoCurRead\circ\UltraCurProj
=
\CurRead\circ\ProtoCurProj\circ\UltraCurProj,
\label{eq:ultra_cur_map}
\end{equation}
\begin{equation}
\UltraHomRead\equiv \ProtoHomRead\circ\UltraHomProj
=
\HomRead\circ\ProtoHomProj\circ\UltraHomProj.
\label{eq:ultra_hom_map}
\end{equation}
The still more primitive microphysical exact-preservation functional is
\begin{equation}
\begin{aligned}
\UltraFunctional[\UltraTensor,\UltraRelabel,\UltraFree]
\equiv\;
&\ProtoFunctional[
\UltraCarProj\UltraTensor,\UltraCurProj\UltraRelabel,\UltraHomProj\UltraFree]
\\
&+\frac12\norm{\TensorDefect\UltraTensor}_{\mathrm{def,car}}^2
+\frac12\norm{\CurDefect\UltraRelabel}_{\mathrm{def,cur}}^2
+\frac12\norm{\HomDefect\UltraFree}_{\mathrm{def,hom}}^2.
\end{aligned}
\label{eq:Fultra}
\end{equation}
\end{definition}

\begin{remark}
At current scope the microphysical functional \eqref{eq:Fultra} is intentionally minimal. The defect-null kernel sectors carry exactly the ambient projected-sector construction already recorded one layer above. The defect-bearing complements are not allowed to create new visible slots or new unsuppressed observer data if exact preservation is to survive. Their only honest current-scope role is therefore to enter through the strictly positive quadratic defect penalties.
\end{remark}

\begin{proposition}[Why the ambient substrate spaces are forced at current scope]
\label{prop:ambient_forced}
If a still more primitive microphysical sector is to preserve the same downstream carrier-origin functional \eqref{eq:car}, the same exact-preservation normal form \eqref{eq:Fdeep}, and the same one operative metric, then its unsuppressed ambient sector is forced to be the defect-null sector
\[
\ProtoTensorClass=\ker\TensorDefect,\qquad
\ProtoRelabelClass=\ker\CurDefect,\qquad
\ProtoFreeClass=\ker\HomDefect.
\]
\end{proposition}

\begin{proof}
Any microphysical direction with nonzero defect changes one of the quantities measured by \(\TensorDefect\), \(\CurDefect\), or \(\HomDefect\). By construction those quantities do not appear as visible slots anywhere in the recorded carrier-origin functional \eqref{eq:car}, the ambient projected-sector functional \eqref{eq:Fproto}, or the exact-preservation normal form \eqref{eq:Fdeep}. If such defect-bearing directions were left unsuppressed, the microphysical theory would either introduce additional visible data beyond the recorded package or alter the stationary reduction away from \eqref{eq:Fproto}. Therefore the only admissible unsuppressed sector is the common defect-null sector, namely \eqref{eq:proto_kernels}. This is exactly the ambient substrate layer that can preserve the recorded downstream package.
\end{proof}

\begin{proposition}[Why the ambient projectors and raw deeper sectors are forced]
\label{prop:ambient_raw}
Once the exact-preserving ambient sectors \eqref{eq:proto_kernels} are fixed, the orthogonal projectors
\[
\ProtoCarProj,\qquad
\ProtoCurProj,\qquad
\ProtoHomProj
\]
onto the raw deeper sectors and the raw deeper spaces
\[
\DeepTensorClass,\qquad
\DeepRelabelClass,\qquad
\DeepFreeClass
\]
remain forced exactly as in the recorded ambient projected-sector construction \eqref{eq:proto_proj}--\eqref{eq:proto_kill}.
\end{proposition}

\begin{proof}
Definition \ref{def:proto} is part of the fixed target data of the present note. It records that, inside each exact-preserving ambient space, only the projected sectors \(\DeepTensorClass=\operatorname{ran}\ProtoCarProj\), \(\DeepRelabelClass=\operatorname{ran}\ProtoCurProj\), and \(\DeepFreeClass=\operatorname{ran}\ProtoHomProj\) are allowed to feed the exact-preserving visible readouts, while the ambient complements are positively gapped spectators. Since the present microphysical layer is required to recover that same ambient functional \eqref{eq:Fproto}, those same ambient projectors and those same raw deeper sectors remain forced without modification once the ambient spaces themselves are recovered.
\end{proof}

\begin{proposition}[Why the codomain-preserving visible maps are forced]
\label{prop:micro_codomains}
At current exact-preservation scope, the only admissible visible maps from the microphysical layer are the factorized maps \eqref{eq:ultra_car_map}--\eqref{eq:ultra_hom_map}. In particular, the only visible codomains remain
\[
\CarrierClass,\qquad
T^*\Sub,\qquad
\CompFree.
\]
\end{proposition}

\begin{proof}
The fixed carrier-origin functional \eqref{eq:car} has exactly three visible slots and no others. Its tensor slot is \(\CarrierSeed\in\CarrierClass\), because the branch-compatible projector \(\CarrierProj\), the source loading, and the later primitive/coarse eliminations are defined there. Its relabeling slot is the covector current \(\RelabelSeed\in T^*\Sub\), because the only current-scope relabeling imprint of the tensor sector enters through the covector \(\SubGaugeVec(\CarrierProj\CarrierSeed)\). Its homogeneous slot is \(\FreeSeed\in\CompFree\), because that is the source-free kernel of \(\SubReducedEin\) and the only homogeneous slot compatible with the later \(\Pi_{\mathrm{free}}H=0\) argument.

No additional visible slot is allowed if the carrier-origin functional is to remain exactly \eqref{eq:car}. Therefore every visible microphysical map must land in those same three codomains. By Proposition \ref{prop:ambient_forced}, exact preservation first requires projection onto the defect-null ambient sectors via \(\UltraCarProj,\UltraCurProj,\UltraHomProj\). By Proposition \ref{prop:ambient_raw}, inside those ambient sectors only the projected raw deeper sectors \(\operatorname{ran}\ProtoCarProj,\operatorname{ran}\ProtoCurProj,\operatorname{ran}\ProtoHomProj\) may feed visible data. Thus the only admissible current-scope visible maps are exactly \eqref{eq:ultra_car_map}--\eqref{eq:ultra_hom_map}.
\end{proof}

\begin{proposition}[Why the hidden coercive sectors and positive gaps are forced]
\label{prop:kernelquadratic}
At current exact-preservation scope, two hidden quadratic sectors are forced.
\begin{enumerate}
    \item The defect-bearing microphysical complements \(\operatorname{ran}\UltraCarPerp\), \(\operatorname{ran}\UltraCurPerp\), and \(\operatorname{ran}\UltraHomPerp\) are coercively gapped with strictly positive gaps \(\UltraCarGap\), \(\UltraCurGap\), and \(\UltraHomGap\).
    \item Inside the recovered ambient projected-sector construction, the hidden readout-null directions \(\ker\CarRead\), \(\ker\CurRead\), and \(\ker\HomRead\) still carry the same positive quadratic sector with strictly positive gaps \(\PreCarGap\), \(\CurrentGap\), and \(\HomGap\).
\end{enumerate}
\end{proposition}

\begin{proof}
Item (1) is exactly the coercive content of \eqref{eq:defect_car}--\eqref{eq:defect_hom}. A defect-bearing microphysical direction is invisible to the recorded downstream package but would remain an unsuppressed exact-breaking degree of freedom if it were not penalized. The defect estimates therefore force strictly positive gaps on the microphysical complements.

Item (2) is built into the fixed ambient exact-preservation functional \eqref{eq:Fproto} and the fixed exact-preservation normal form \eqref{eq:Fdeep}. Once the ambient projected-sector construction is recovered, the ambient complements remain the positively gapped spectator sectors with gaps \(\ProtoCarGap,\ProtoCurGap,\ProtoHomGap\), and inside the raw deeper sectors the readout-null directions remain the positively gapped kernel sectors with gaps \(\PreCarGap,\CurrentGap,\HomGap\). Thus the kernel-coercive quadratic sector of the exact-preservation normal form is not lost but recovered unchanged.
\end{proof}

\section{Exact Recovery of the Same Ambient Exact-Preservation Sector}

\begin{definition}[Microphysical coarse-graining to the ambient sector]
\label{def:Cultra}
For fixed ambient data \((\ProtoTensor,\ProtoRelabel,\ProtoFree)\), define
\begin{equation}
\begin{aligned}
\mathscr C_{\mathrm{micro}}
[\ProtoTensor,\ProtoRelabel,\ProtoFree;
\UltraTensor,\UltraRelabel,\UltraFree,
\Lambda_{\mathrm{car}},\Lambda_{\mathrm{cur}},\Lambda_{\mathrm{hom}}]
\equiv\;
&\UltraFunctional[\UltraTensor,\UltraRelabel,\UltraFree]
\\
&+\left\langle
\Lambda_{\mathrm{car}},
\ProtoTensor-\UltraCarProj\UltraTensor
\right\rangle_{\mathrm{proto,car}}
\\
&+\left\langle
\Lambda_{\mathrm{cur}},
\ProtoRelabel-\UltraCurProj\UltraRelabel
\right\rangle_{\mathrm{proto,cur}}
\\
&+\left\langle
\Lambda_{\mathrm{hom}},
\ProtoFree-\UltraHomProj\UltraFree
\right\rangle_{\mathrm{proto,hom}},
\end{aligned}
\label{eq:Cultra}
\end{equation}
where
\[
\Lambda_{\mathrm{car}}\in\ProtoTensorClass,\qquad
\Lambda_{\mathrm{cur}}\in\ProtoRelabelClass,\qquad
\Lambda_{\mathrm{hom}}\in\ProtoFreeClass.
\]
The microphysical stationary-value functional is
\begin{equation}
\UltraToProto[\ProtoTensor,\ProtoRelabel,\ProtoFree]
\equiv
\operatorname*{stat}_{\UltraTensor,\UltraRelabel,\UltraFree,\Lambda_{\mathrm{car}},\Lambda_{\mathrm{cur}},\Lambda_{\mathrm{hom}}}
\mathscr C_{\mathrm{micro}}.
\label{eq:Fultrastat}
\end{equation}
\end{definition}

\begin{theorem}[Exact recovery of the same ambient exact-preservation functional]
\label{thm:recoverproto}
The stationary equations of \eqref{eq:Cultra} enforce
\begin{equation}
\UltraCarProj\UltraTensor=\ProtoTensor,
\qquad
\UltraCurProj\UltraRelabel=\ProtoRelabel,
\qquad
\UltraHomProj\UltraFree=\ProtoFree,
\label{eq:ultra_constraints_1}
\end{equation}
together with
\begin{equation}
\UltraCarPerp\UltraTensor=0,
\qquad
\UltraCurPerp\UltraRelabel=0,
\qquad
\UltraHomPerp\UltraFree=0.
\label{eq:ultra_constraints_2}
\end{equation}
Consequently,
\begin{equation}
\UltraToProto[\ProtoTensor,\ProtoRelabel,\ProtoFree]
=
\ProtoFunctional[\ProtoTensor,\ProtoRelabel,\ProtoFree].
\label{eq:ultra_to_proto}
\end{equation}
\end{theorem}

\begin{proof}
Stationarity of \eqref{eq:Cultra} with respect to the multipliers \(\Lambda_{\mathrm{car}},\Lambda_{\mathrm{cur}},\Lambda_{\mathrm{hom}}\) gives \eqref{eq:ultra_constraints_1}. The complementary microphysical directions \(\UltraCarPerp\UltraTensor\), \(\UltraCurPerp\UltraRelabel\), and \(\UltraHomPerp\UltraFree\) do not appear in the multiplier constraints because the projectors annihilate them. They appear in \eqref{eq:Cultra} only through the strictly positive quadratic defect penalties in \eqref{eq:Fultra}. Variation with respect to those complementary directions gives
\[
\TensorDefect^*\TensorDefect\,\UltraCarPerp\UltraTensor=0,
\qquad
\CurDefect^*\CurDefect\,\UltraCurPerp\UltraRelabel=0,
\qquad
\HomDefect^*\HomDefect\,\UltraHomPerp\UltraFree=0.
\]
Using the coercive defect bounds \eqref{eq:defect_car}--\eqref{eq:defect_hom}, one obtains \eqref{eq:ultra_constraints_2}. Substituting \eqref{eq:ultra_constraints_1}--\eqref{eq:ultra_constraints_2} into \eqref{eq:Cultra} removes the multiplier terms and annihilates the microphysical defect penalties. What remains is exactly
\[
\ProtoFunctional[\ProtoTensor,\ProtoRelabel,\ProtoFree].
\]
This proves \eqref{eq:ultra_to_proto}.
\end{proof}

\begin{corollary}[Same ambient projected-sector package]
\label{cor:sameproto}
The still deeper microphysical sector reproduces exactly the same ambient substrate spaces \(\ProtoTensorClass,\ProtoRelabelClass,\ProtoFreeClass\), the same ambient projectors \(\ProtoCarProj,\ProtoCurProj,\ProtoHomProj\), the same factorized visible maps \(\ProtoCarRead,\ProtoCurRead,\ProtoHomRead\), the same positive ambient spectator gaps \(\ProtoCarGap,\ProtoCurGap,\ProtoHomGap\), and the same ambient exact-preservation functional \eqref{eq:Fproto} recorded in the previous note.
\end{corollary}

\begin{proof}
This is exactly Theorem \ref{thm:recoverproto} together with Propositions \ref{prop:ambient_forced}, \ref{prop:ambient_raw}, \ref{prop:micro_codomains}, and \ref{prop:kernelquadratic}.
\end{proof}

\section{Exact Recovery of the Same Exact-Preservation Normal Form}

\begin{definition}[Ambient coarse-graining to the raw deeper sector]
\label{def:Cproto}
For fixed raw deeper data \((\DeepTensor,\DeepRelabel,\DeepFree)\), define
\begin{equation}
\begin{aligned}
\mathscr C_{\mathrm{proto}}
[\DeepTensor,\DeepRelabel,\DeepFree;
\ProtoTensor,\ProtoRelabel,\ProtoFree,
\Lambda_{\mathrm{car}},\Lambda_{\mathrm{cur}},\Lambda_{\mathrm{hom}}]
\equiv\;
&\ProtoFunctional[\ProtoTensor,\ProtoRelabel,\ProtoFree]
\\
&+\left\langle
\Lambda_{\mathrm{car}},
\DeepTensor-\ProtoCarProj\ProtoTensor
\right\rangle_{\mathrm{deep,car}}
\\
&+\left\langle
\Lambda_{\mathrm{cur}},
\DeepRelabel-\ProtoCurProj\ProtoRelabel
\right\rangle_{\mathrm{deep,cur}}
\\
&+\left\langle
\Lambda_{\mathrm{hom}},
\DeepFree-\ProtoHomProj\ProtoFree
\right\rangle_{\mathrm{deep,hom}},
\end{aligned}
\label{eq:Cproto}
\end{equation}
where
\[
\Lambda_{\mathrm{car}}\in\DeepTensorClass,\qquad
\Lambda_{\mathrm{cur}}\in\DeepRelabelClass,\qquad
\Lambda_{\mathrm{hom}}\in\DeepFreeClass.
\]
The ambient stationary-value functional is
\begin{equation}
\ProtoToDeep[\DeepTensor,\DeepRelabel,\DeepFree]
\equiv
\operatorname*{stat}_{\ProtoTensor,\ProtoRelabel,\ProtoFree,\Lambda_{\mathrm{car}},\Lambda_{\mathrm{cur}},\Lambda_{\mathrm{hom}}}
\mathscr C_{\mathrm{proto}}.
\label{eq:Fprotostat}
\end{equation}
\end{definition}

\begin{theorem}[Exact recovery of the same exact-preservation normal form]
\label{thm:recoverdeep}
The stationary equations of \eqref{eq:Cproto} enforce
\begin{equation}
\ProtoCarProj\ProtoTensor=\DeepTensor,
\qquad
\ProtoCurProj\ProtoRelabel=\DeepRelabel,
\qquad
\ProtoHomProj\ProtoFree=\DeepFree,
\label{eq:proto_constraints_1}
\end{equation}
together with
\begin{equation}
\ProtoCarPerp\ProtoTensor=0,
\qquad
\ProtoCurPerp\ProtoRelabel=0,
\qquad
\ProtoHomPerp\ProtoFree=0.
\label{eq:proto_constraints_2}
\end{equation}
Consequently,
\begin{equation}
\ProtoToDeep[\DeepTensor,\DeepRelabel,\DeepFree]
=
\DeepFunctional[\DeepTensor,\DeepRelabel,\DeepFree].
\label{eq:proto_to_deep}
\end{equation}
\end{theorem}

\begin{proof}
Stationarity of \eqref{eq:Cproto} with respect to the multipliers \(\Lambda_{\mathrm{car}},\Lambda_{\mathrm{cur}},\Lambda_{\mathrm{hom}}\) gives \eqref{eq:proto_constraints_1}. The complementary ambient directions \(\ProtoCarPerp\ProtoTensor\), \(\ProtoCurPerp\ProtoRelabel\), and \(\ProtoHomPerp\ProtoFree\) do not appear in the multiplier constraints because the projectors annihilate them. They appear in \eqref{eq:Cproto} only through the strictly positive quadratic penalties in \eqref{eq:Fproto}. Therefore variation with respect to those complementary ambient directions gives
\[
\ProtoCarGap\,\ProtoCarPerp\ProtoTensor=0,
\qquad
\ProtoCurGap\,\ProtoCurPerp\ProtoRelabel=0,
\qquad
\ProtoHomGap\,\ProtoHomPerp\ProtoFree=0,
\]
and since \(\ProtoCarGap,\ProtoCurGap,\ProtoHomGap>0\), equation \eqref{eq:proto_constraints_2} follows.

Substituting \eqref{eq:proto_constraints_1}--\eqref{eq:proto_constraints_2} into \eqref{eq:Cproto} removes the multiplier terms and annihilates the ambient spectator penalties. What remains is exactly
\[
\DeepFunctional[\DeepTensor,\DeepRelabel,\DeepFree].
\]
This proves \eqref{eq:proto_to_deep}.
\end{proof}

\begin{corollary}[Same raw deeper readout package]
\label{cor:samedeep}
The still deeper microphysical sector reproduces exactly the same raw deeper readout spaces \(\DeepTensorClass,\DeepRelabelClass,\DeepFreeClass\), the same codomain-preserving maps \(\CarRead,\CurRead,\HomRead\), the same hidden kernel gaps \(\PreCarGap,\CurrentGap,\HomGap\), and the same exact-preservation normal-form functional \eqref{eq:Fdeep} recorded one layer above.
\end{corollary}

\begin{proof}
This is exactly Corollary \ref{cor:sameproto}, Theorem \ref{thm:recoverdeep}, and Proposition \ref{prop:kernelquadratic}.
\end{proof}

\begin{corollary}[Same precursor-selection normal form]
\label{cor:precursoravailable}
Because the exact-preservation normal-form functional \eqref{eq:Fdeep} is unchanged, the same precursor-selection normal form of the previous note remains available without modification.
\end{corollary}

\begin{proof}
The previous note's precursor-selection construction depends only on the data
\[
\DeepTensorClass,\qquad
\DeepRelabelClass,\qquad
\DeepFreeClass,
\]
the maps \(\CarRead,\CurRead,\HomRead\), and the positive hidden-sector gaps appearing in \eqref{eq:Fdeep}. By Corollary \ref{cor:samedeep}, those data are reproduced exactly.
\end{proof}

\section{Recovery of the Same Carrier-Origin Lift}

\begin{definition}[Raw deeper coarse-graining functional]
\label{def:deep_cg}
For fixed carrier-origin readouts \((\CarrierSeed,\RelabelSeed,\FreeSeed)\), define
\begin{equation}
\begin{aligned}
\mathscr C_{\mathrm{deep}}
[\CarrierSeed,\RelabelSeed,\FreeSeed;
\DeepTensor,\DeepRelabel,\DeepFree,
\StrainMult,\NeutralMult,\FreeMult]
\equiv\;
&\DeepFunctional[\DeepTensor,\DeepRelabel,\DeepFree]
\\
&+\int_{\Sub} d^4X\,\sqrt{G}\,
\StrainMult^{AB}\bigl(\CarrierSeed-\CarRead\DeepTensor\bigr)_{AB}
\\
&+\int_{\Sub} d^4X\,\sqrt{G}\,
\NeutralMult^A\bigl(\RelabelSeed-\CurRead\DeepRelabel\bigr)_A
\\
&+\int_{\Sub} d^4X\,\sqrt{G}\,
\FreeMult^{AB}\bigl(\FreeSeed-\HomRead\DeepFree\bigr)_{AB},
\end{aligned}
\label{eq:Cdeep}
\end{equation}
where \(\StrainMult\in\CarrierClass\), \(\NeutralMult\in T^*\Sub\), and \(\FreeMult\in\CompFree\). The stationary-value functional is
\begin{equation}
\DeepToCar[\CarrierSeed,\RelabelSeed,\FreeSeed]
\equiv
\operatorname*{stat}_{\DeepTensor,\DeepRelabel,\DeepFree,\StrainMult,\NeutralMult,\FreeMult}
\mathscr C_{\mathrm{deep}}.
\label{eq:Fdeepstat}
\end{equation}
\end{definition}

\begin{theorem}[Exact recovery of the same carrier-origin functional]
\label{thm:recover_car}
The stationary equations of \eqref{eq:Cdeep} enforce
\begin{equation}
\CarRead\DeepTensor=\CarrierSeed,
\qquad
\CurRead\DeepRelabel=\RelabelSeed,
\qquad
\HomRead\DeepFree=\FreeSeed,
\label{eq:deep_constraints_1}
\end{equation}
together with
\begin{equation}
\DeepTensor^{\perp}=0,
\qquad
\DeepRelabel^{\perp}=0,
\qquad
\DeepFree^{\perp}=0.
\label{eq:deep_constraints_2}
\end{equation}
Consequently,
\begin{equation}
\DeepToCar[\CarrierSeed,\RelabelSeed,\FreeSeed]
=
\CarrierFunctional[\CarrierSeed,\RelabelSeed,\FreeSeed].
\label{eq:deep_to_carrier}
\end{equation}
\end{theorem}

\begin{proof}
By Theorem \ref{thm:recoverdeep}, the deeper functional appearing in \eqref{eq:Cdeep} is exactly the same exact-preservation normal form \eqref{eq:Fdeep} recorded in the previous note. Stationarity of \eqref{eq:Cdeep} with respect to the multipliers \(\StrainMult\), \(\NeutralMult\), and \(\FreeMult\) gives \eqref{eq:deep_constraints_1}. The readout maps \(\CarRead\), \(\CurRead\), and \(\HomRead\) vanish on their kernels, so the hidden pieces \(\DeepTensor^{\perp}\), \(\DeepRelabel^{\perp}\), and \(\DeepFree^{\perp}\) appear only through the positive quadratic penalties in \eqref{eq:Fdeep}. Variation therefore gives
\[
\PreCarGap\,\DeepTensor^{\perp}=0,
\qquad
\CurrentGap\,\DeepRelabel^{\perp}=0,
\qquad
\HomGap\,\DeepFree^{\perp}=0,
\]
and because the coefficients are strictly positive, \eqref{eq:deep_constraints_2} follows. Substituting \eqref{eq:deep_constraints_1}--\eqref{eq:deep_constraints_2} into \eqref{eq:Cdeep} removes the multiplier terms and annihilates the positive hidden-sector penalties, leaving exactly
\[
\CarrierFunctional[\CarrierSeed,\RelabelSeed,\FreeSeed].
\]
This proves \eqref{eq:deep_to_carrier}.
\end{proof}

\begin{corollary}[Same carrier-origin readout package]
\label{cor:same_car_package}
The still deeper microphysical sector reproduces exactly the same carrier-origin tensor seed \(\CarrierSeed\in\CarrierClass\), the same relabeling current \(\RelabelSeed\in T^*\Sub\), the same free homogeneous seed \(\FreeSeed\in\CompFree\), and the same carrier-origin functional \eqref{eq:car}.
\end{corollary}

\begin{proof}
This is exactly Theorem \ref{thm:recover_car}.
\end{proof}

\section{Recovery of the Same Primitive Reservoir and Microscopic Reservoir}

\begin{definition}[Fixed carrier coarse-graining functional]
\label{def:car_cg}
For fixed primitive readouts \((\PrimStrain,\PrimNeutral,\PrimFree)\), define
\begin{equation}
\begin{aligned}
\mathscr C_{\mathrm{car}}
[\PrimStrain,\PrimNeutral,\PrimFree;
\CarrierSeed,\RelabelSeed,\FreeSeed,
\StrainMult,\NeutralMult,\FreeMult]
\equiv\;
&\CarrierFunctional[\CarrierSeed,\RelabelSeed,\FreeSeed]
\\
&+\int_{\Sub} d^4X\,\sqrt{G}\,
\StrainMult^{AB}\bigl(\PrimStrain-\CarrierProj\CarrierSeed\bigr)_{AB}
\\
&+\int_{\Sub} d^4X\,\sqrt{G}\,
\NeutralMult^A(\PrimNeutral-\RelabelSeed)_A
\\
&+\int_{\Sub} d^4X\,\sqrt{G}\,
\FreeMult^{AB}(\PrimFree-\FreeSeed)_{AB}.
\end{aligned}
\label{eq:Ccar}
\end{equation}
Its stationary-value functional is
\begin{equation}
\CarrierToPrim[\PrimStrain,\PrimNeutral,\PrimFree]
\equiv
\operatorname*{stat}_{\CarrierSeed,\RelabelSeed,\FreeSeed,\StrainMult,\NeutralMult,\FreeMult}
\mathscr C_{\mathrm{car}}.
\label{eq:Fcarstat}
\end{equation}
\end{definition}

\begin{theorem}[Exact recovery of the same primitive reservoir]
\label{thm:recoverprimitive}
The still deeper microphysical sector reproduces exactly the same primitive reservoir:
\begin{equation}
\CarrierToPrim[\PrimStrain,\PrimNeutral,\PrimFree]
=
\PrimFunctional[\PrimStrain,\PrimNeutral,\PrimFree].
\label{eq:primitive_exact}
\end{equation}
\end{theorem}

\begin{proof}
By Theorem \ref{thm:recover_car}, the carrier functional appearing in \eqref{eq:Ccar} is exactly the same carrier-origin functional \eqref{eq:car} used in the primitive-reservoir carrier-origin note. The stationarity equations of \eqref{eq:Ccar} therefore again enforce
\[
\CarrierProj\CarrierSeed=\PrimStrain,
\qquad
\RelabelSeed=\PrimNeutral,
\qquad
\FreeSeed=\PrimFree,
\qquad
\CarrierPerp\CarrierSeed=0.
\]
Substituting these relations into \eqref{eq:Ccar} removes the multiplier terms and leaves exactly \eqref{eq:primitive}. This proves \eqref{eq:primitive_exact}.
\end{proof}

\begin{corollary}[Same reduced carrier Hessian]
\label{cor:hessian}
Eliminating the neutral current from the recovered primitive reservoir gives the same reduced primitive functional
\begin{equation}
\begin{aligned}
\int_{\Sub} d^4X\,\sqrt{G}\,
\Bigl[
\frac12 \PrimStrain^{AB}\SubReducedEin(\PrimStrain)_{AB}
+\frac{\FreeStiff}{2}\PrimFree^{AB}\PrimFree_{AB}
+(\ResidualLift)^{AB}\PrimStrain_{AB}
\Bigr],
\end{aligned}
\label{eq:Fredprimitive}
\end{equation}
whose quadratic \(\PrimStrain\)-Hessian is the same completion kernel
\begin{equation}
\CompKernel(H_1,H_2)
\equiv
\int_{\Sub} d^4X\,\sqrt{G}\,
H_1^{AB}\SubReducedEin(H_2)_{AB}.
\label{eq:kernel}
\end{equation}
\end{corollary}

\begin{proof}
By \eqref{eq:primitive_exact}, variation with respect to \(\PrimNeutral_A\) gives
\[
\PrimNeutral_A=\SubGaugeVec(\PrimStrain)_A.
\]
Substituting this back into \eqref{eq:primitive} yields \eqref{eq:Fredprimitive}. Bilinearization of the quadratic \(\frac12\int \PrimStrain^{AB}\SubReducedEin(\PrimStrain)_{AB}\) term gives \eqref{eq:kernel}.
\end{proof}

\begin{definition}[Fixed primitive coarse-graining functional]
\label{def:prim_cg}
For fixed coarse completion data \((H,\GaugeAux)\in\CompClass\times T^*\Sub\), define
\begin{equation}
\begin{aligned}
\mathscr C_{\mathrm{prim}}
[H,\GaugeAux;
\PrimStrain,\PrimNeutral,\PrimFree,
\StrainMult,\NeutralMult,\FreeMult]
\equiv\;
&\PrimFunctional[\PrimStrain,\PrimNeutral,\PrimFree]
\\
&+\int_{\Sub} d^4X\,\sqrt{G}\,
\StrainMult^{AB}(H-\PrimStrain)_{AB}
\\
&+\int_{\Sub} d^4X\,\sqrt{G}\,
\NeutralMult^A(\GaugeAux-\PrimNeutral)_A
\\
&+\int_{\Sub} d^4X\,\sqrt{G}\,
\FreeMult^{AB}\bigl((\FreeProj H)-\PrimFree\bigr)_{AB}.
\end{aligned}
\label{eq:Cprim}
\end{equation}
\end{definition}

\begin{theorem}[Exact recovery of the same microscopic-equilibrium reservoir]
\label{thm:recovermic}
Using \eqref{eq:primitive_exact} inside \eqref{eq:Cprim}, the same stationary-value macroscopic reservoir \eqref{eq:mic} is recovered:
\begin{equation}
\operatorname*{stat}_{\PrimStrain,\PrimNeutral,\PrimFree,\StrainMult,\NeutralMult,\FreeMult}
\mathscr C_{\mathrm{prim}}
=
\MicFunctional[H,\GaugeAux].
\label{eq:micrecover}
\end{equation}
\end{theorem}

\begin{proof}
By Theorem \ref{thm:recoverprimitive}, the primitive functional appearing in \eqref{eq:Cprim} is exactly the same \(\PrimFunctional\) already used in the equilibrium-reservoir primitive-origin note. The stationarity equations of \eqref{eq:Cprim} therefore again enforce
\[
\PrimStrain=H,
\qquad
\PrimNeutral=\GaugeAux,
\qquad
\PrimFree=\FreeProj H,
\]
and substitution yields \eqref{eq:mic}. Hence \eqref{eq:micrecover} follows.
\end{proof}

\section{Recovery of the Same Completion Solve and One Operative Metric}

\begin{theorem}[Same stationary completion equations]
\label{thm:stationary}
Let \((\CompLift,\GaugeAux)\) be a stationary point of the recovered macroscopic reservoir \(\MicFunctional\). Then
\begin{equation}
\GaugeAux_A=\SubGaugeVec(\CompLift)_A,
\label{eq:Yeq}
\end{equation}
and
\begin{equation}
\SubReducedEin(\CompLift)_{AB}
+\FreeStiff(\FreeProj\CompLift)_{AB}
=
-\ResidualLift_{AB}.
\label{eq:Heq}
\end{equation}
\end{theorem}

\begin{proof}
Theorem \ref{thm:recovermic} shows that the macroscopic reservoir is exactly \eqref{eq:mic}. Variation of \eqref{eq:mic} with respect to \(\GaugeAux_A\) gives \eqref{eq:Yeq}. Substituting \eqref{eq:Yeq} back into \eqref{eq:mic} yields the reduced functional
\[
\int_{\Sub} d^4X\,\sqrt{G}\,
\Bigl[
\frac12 H^{AB}\SubReducedEin(H)_{AB}
+\frac{\FreeStiff}{2}(\FreeProj H)^{AB}(\FreeProj H)_{AB}
+(\ResidualLift)^{AB}H_{AB}
\Bigr],
\]
and variation with respect to \(H_{AB}\) gives \eqref{eq:Heq}.
\end{proof}

\begin{proposition}[Same zero-free-mode outcome]
\label{prop:freezero}
The stationary completion tensor satisfies
\begin{equation}
\FreeProj\CompLift=0.
\label{eq:freezero}
\end{equation}
\end{proposition}

\begin{proof}
Apply \(\FreeProj\) to \eqref{eq:Heq}. Since \(\CompFree=\ker\SubReducedEin\), one has \(\FreeProj\SubReducedEin(\CompLift)=0\). By \eqref{eq:sourcefree}, \(\FreeProj\ResidualLift=0\). Therefore
\[
\FreeStiff\,\FreeProj\CompLift=0.
\]
Because \(\FreeStiff>0\), equation \eqref{eq:freezero} follows.
\end{proof}

\begin{corollary}[Same substrate completion solve]
\label{cor:subsolve}
The stationary completion tensor satisfies
\begin{equation}
\SubReducedEin(\CompLift)_{AB}
=
-\ResidualLift_{AB},
\qquad
\FreeProj\CompLift=0.
\label{eq:subsolve}
\end{equation}
\end{corollary}

\begin{proof}
Substituting \eqref{eq:freezero} into \eqref{eq:Heq} gives \eqref{eq:subsolve}.
\end{proof}

\begin{definition}[Completed bridge and observer metrics]
\label{def:completed}
Define
\begin{equation}
\CompShift_{\mu\nu}
\equiv
\ObserverMap^*\CompLift_{AB},
\qquad
\CompletedMetric
\equiv
\ObserverMetric+\BaseShift+\ResponseShift+\CompShift.
\label{eq:completedmetric}
\end{equation}
\end{definition}

\begin{theorem}[Same observer solve and one operative metric]
\label{thm:observer}
The still deeper microphysical sector reproduces exactly the same observer completion solve
\begin{equation}
\ReducedEin(\CompShift)_{\mu\nu}
=
-\ResidualCore{}_{\mu\nu}.
\label{eq:observereq}
\end{equation}
Consequently the same completed observer metric \(\CompletedMetric\) remains the unique operative metric of the branch, with the same universal matter route and the same observer-equation removal of the former genuine quartic residual sector.
\end{theorem}

\begin{proof}
Apply \(\ObserverMap^*\) to \eqref{eq:subsolve}. Using \eqref{eq:intertwine2} and \(\ObserverMap^*\ResidualLift=\ResidualCore\), one obtains \eqref{eq:observereq}. No new observer metric, bridge map, or matter-coupling rule is introduced anywhere in the microphysical sector, in the recovered ambient projected-sector layer, or in the recovered raw deeper sector. Therefore the same completed metric \(\CompletedMetric\) remains the unique operative metric, exactly as in the previous completion notes.
\end{proof}

\begin{corollary}[Recorded gain package is preserved]
\label{cor:gains}
Because the recovered exact-preservation normal form \eqref{eq:Fdeep}, carrier-origin functional \eqref{eq:car}, primitive reservoir \eqref{eq:primitive}, macroscopic reservoir \eqref{eq:mic}, and stationary equations \eqref{eq:Yeq}--\eqref{eq:observereq} are unchanged, the recorded Phase D / E infrared-spectrum and universal-coupling verdicts, the recorded Phase F controlled GR limit, and the zero-free-mode verdict \(\Pi_{\mathrm{free}}H=0\) are preserved without modification.
\end{corollary}

\begin{proof}
The present note changes only the still deeper origin of the ambient projected-sector construction and hence of the exact-preservation normal form. It leaves the ambient functional, the deeper normal form, the carrier-origin functional, the primitive reservoir, the macroscopic completion functional, the completion solve, and the one operative metric unchanged. Therefore all previously recorded observer-level consequences of that same completion package survive verbatim.
\end{proof}

\section{Claim-Boundary Verdict}

\begin{theorem}[Exact-preservation normal-form microphysical-origin verdict]
\label{thm:verdict}
At the disciplined scope of the present note, the positive result is exactly the following.
\begin{enumerate}
    \item The ambient projected-sector construction of the previous note is no longer left as the present deepest disciplined block.
    \item The ambient substrate spaces \(\ProtoTensorClass\), \(\ProtoRelabelClass\), and \(\ProtoFreeClass\) are recovered as the defect-null exact-preserving sectors of still more primitive microphysical spaces \(\UltraTensorClass\), \(\UltraRelabelClass\), and \(\UltraFreeClass\).
    \item The orthogonal kernel projectors \(\UltraCarProj\), \(\UltraCurProj\), and \(\UltraHomProj\) are forced by that defect-null structure, and inside the recovered ambient layer the same ambient projectors \(\ProtoCarProj\), \(\ProtoCurProj\), and \(\ProtoHomProj\) again recover the same raw deeper sectors \(\DeepTensorClass\), \(\DeepRelabelClass\), and \(\DeepFreeClass\).
    \item The visible microphysical maps are forced to factor through the recovered ambient and raw-sector projectors into the same codomain-preserving maps landing in \(\CarrierClass\), \(T^*\Sub\), and \(\CompFree\), because those remain the only visible slots of the fixed carrier-origin functional.
    \item The defect-bearing microphysical complements carry the strictly positive gaps \(\UltraCarGap\), \(\UltraCurGap\), and \(\UltraHomGap\); the ambient complements carry the strictly positive gaps \(\ProtoCarGap\), \(\ProtoCurGap\), and \(\ProtoHomGap\); and inside the raw deeper sectors the hidden readout-null directions \(\ker\CarRead\), \(\ker\CurRead\), and \(\ker\HomRead\) remain the kernel-coercive positively gapped sectors with gaps \(\PreCarGap\), \(\CurrentGap\), and \(\HomGap\).
    \item Eliminating the microphysical defect layer reproduces exactly the same ambient exact-preservation functional \(\ProtoFunctional[\ProtoTensor,\ProtoRelabel,\ProtoFree]\).
    \item Eliminating the recovered ambient layer reproduces exactly the same exact-preservation normal-form functional \(\DeepFunctional[\DeepTensor,\DeepRelabel,\DeepFree]\).
    \item The already recorded downstream eliminations therefore reproduce exactly the same carrier-origin functional, the same primitive reservoir, the same microscopic-equilibrium reservoir, the same \(\Pi_{\mathrm{free}}H=0\) outcome, the same substrate solve \(\SubReducedEin(\CompLift)=-\ResidualLift\), the same observer solve \(\ReducedEin(\CompShift)=-\ResidualCore\), the same completion solve \(\Sigma^{\mathrm{cmp}}\), and the same one operative metric.
    \item Stronger promotion language is still not justified here, because the present note supplies only the microphysical origin of the ambient projected-sector construction at current scope, not a final ultraviolet completion or a theorem wider than the current one-metric claim boundary.
\end{enumerate}
\end{theorem}

\begin{proof}
Item (1) is the main framing claim of the introduction. Item (2) is Proposition \ref{prop:ambient_forced}. Item (3) is Proposition \ref{prop:ambient_raw} together with Definition \ref{def:ultra}. Item (4) is Proposition \ref{prop:micro_codomains}. Item (5) is Proposition \ref{prop:kernelquadratic}. Item (6) is Theorem \ref{thm:recoverproto}. Item (7) is Theorem \ref{thm:recoverdeep}. Item (8) is Theorem \ref{thm:recover_car}, Theorem \ref{thm:recoverprimitive}, Theorem \ref{thm:recovermic}, Proposition \ref{prop:freezero}, Corollary \ref{cor:subsolve}, Theorem \ref{thm:observer}, and Corollary \ref{cor:gains}. Item (9) is the claim-boundary discipline stated in the introduction.
\end{proof}

\begin{corollary}[What remains next]
The primary active burden moves one layer deeper again. It is no longer to justify why the ambient substrate spaces, their ambient projectors onto the raw deeper sectors, their factorized codomain-preserving maps, their positive spectator gaps, and the same exact-preservation normal form are admissible at current scope; that step is now recorded. The next honest burden is either a still deeper origin or sharper selection principle below the present defect-null microphysical construction itself, or another comparably concrete observer-equation gain on the same one-metric line. The anchored positive-pair continuation remains side work unless it directly supplies one of those.
\end{corollary}

\section{Conclusion}

The positive result of this note is that the ambient projected-sector construction is no longer left unexplained at present scope. It is recovered as the defect-null exact-preserving sector of a still more primitive microphysical layer.

At that deeper level, one begins not with the ambient substrate spaces themselves, but with microphysical spaces
\[
\UltraTensorClass,\qquad
\UltraRelabelClass,\qquad
\UltraFreeClass
\]
equipped with closed defect operators
\[
\TensorDefect,\qquad
\CurDefect,\qquad
\HomDefect.
\]
The admissible exact-preserving ambient sectors are exactly their kernels:
\[
\ProtoTensorClass=\ker\TensorDefect,
\qquad
\ProtoRelabelClass=\ker\CurDefect,
\qquad
\ProtoFreeClass=\ker\HomDefect.
\]
Because these kernels are closed Hilbert subspaces, the orthogonal kernel projectors
\[
\UltraCarProj,\qquad
\UltraCurProj,\qquad
\UltraHomProj
\]
are induced structure rather than a same-layer postulate. The defect-bearing complements are not allowed to become new visible degrees of freedom if the recorded carrier-origin functional is to remain unchanged, so they appear only through strictly positive quadratic penalties:
\[
\UltraCarGap>0,\qquad
\UltraCurGap>0,\qquad
\UltraHomGap>0.
\]

Inside the recovered ambient sectors, the same ambient projectors
\[
\ProtoCarProj,\qquad
\ProtoCurProj,\qquad
\ProtoHomProj
\]
again isolate the same raw deeper sectors
\[
\DeepTensorClass=\mathrm{ran}\,\ProtoCarProj,
\qquad
\DeepRelabelClass=\mathrm{ran}\,\ProtoCurProj,
\qquad
\DeepFreeClass=\mathrm{ran}\,\ProtoHomProj.
\]
The visible microphysical maps are forced to factor first through the microphysical kernel projectors and then through the ambient projectors:
\[
\UltraCarRead=\CarRead\circ\ProtoCarProj\circ\UltraCarProj,
\qquad
\UltraCurRead=\CurRead\circ\ProtoCurProj\circ\UltraCurProj,
\qquad
\UltraHomRead=\HomRead\circ\ProtoHomProj\circ\UltraHomProj.
\]
That factorization is forced because the fixed carrier-origin functional still has exactly three visible slots and no others:
\[
\CarrierClass,\qquad
T^*\Sub,\qquad
\CompFree.
\]
Inside the raw deeper sectors, the hidden readout-null directions are still
\[
\ker\CarRead,\qquad
\ker\CurRead,\qquad
\ker\HomRead,
\]
and they still enter through the same positively gapped kernel-coercive quadratic block:
\[
\PreCarGap>0,\qquad
\CurrentGap>0,\qquad
\HomGap>0.
\]

Constrained elimination of the defect-bearing microphysical layer reproduces exactly the same ambient exact-preservation functional
\[
\ProtoFunctional[\ProtoTensor,\ProtoRelabel,\ProtoFree].
\]
Eliminating the recovered ambient layer then reproduces exactly the same exact-preservation normal-form functional
\[
\DeepFunctional[\DeepTensor,\DeepRelabel,\DeepFree],
\]
and the previously recorded downstream eliminations therefore reproduce exactly the same carrier-origin functional, the same primitive branch-strain / neutral-current / free-mode reservoir, and the same microscopic-equilibrium completion reservoir:
\[
\CarrierFunctional[\CarrierSeed,\RelabelSeed,\FreeSeed],
\qquad
\PrimFunctional[\PrimStrain,\PrimNeutral,\PrimFree],
\qquad
\MicFunctional[H,\GaugeAux].
\]
Hence the same zero-free-mode verdict, the same substrate solve, the same observer solve, the same completion \(\Sigma^{\mathrm{cmp}}\), and the same one operative metric follow once more:
\[
\FreeProj\CompLift=0,
\qquad
\SubReducedEin(\CompLift)=-\ResidualLift,
\qquad
\ReducedEin(\CompShift)=-\ResidualCore.
\]
The recorded Phase D / E and Phase F gains therefore remain intact.

The scope remains honest. This note does not claim a final ultraviolet completion and does not widen the current theorem boundary. Its exact positive result is narrower and precisely the one now needed: the ambient projected-sector construction is no longer primitive at current scope, but is recovered as the defect-null exact-preserving sector of a still more primitive microphysical layer while preserving the same carrier-origin functional, the same primitive reservoir, the same microscopic-equilibrium reservoir, the same \(\Pi_{\mathrm{free}}H=0\) verdict, the same \(\Sigma^{\mathrm{cmp}}\) solve, and the same one operative metric. The next honest burden is one layer deeper again: whether that defect-null microphysical construction can itself be derived from a sharper microscopic symmetry or selection principle, or else superseded by another equally concrete observer-equation gain on the same one-metric line.

\end{document}
